\chapter{网络栈与 WiFi 驱动分层}
\label{chap:netstack}

本章说明 miniOS 的网络代码按 Linux 驱动分层重新组织后的结构：每一层放在哪个目录、对应 Linux 的哪个模块、层与层之间通过什么接口说话，以及一个数据包从 UDP 走到 XR829 固件、再从空口回到 socket 的完整路径。

XR829 的 SDIO 枚举、HIF 寄存器、固件下载等硬件细节见第~\ref{chap:wifi} 章；\file{/boot}（FAT32）和 tftp 写文件见第~\ref{chap:bootfat32} 章。

\section{分层总览}

\begin{txtcode}[basicstyle=\ttfamily\scriptsize]
 用户态  ifconfig / ping / tftp / wifi        (user-t113/*.c)
            │ T113_SYS_NETCTL / T113_SYS_TFTP
 ─────────── src/kernel/syscall.c ─────────────────────────────────────
            │
 net/tftp   tftp.c (协议状态机)  tftp_udp.c (收进内存 → 写 /boot)
            │ udp_open/sendto/recvfrom
 net/ipv4   udp.c  icmp.c  ip_input.c  ip_output.c  arp.c
            devinet.c (地址)  ipconfig.c (DHCP)  af_inet.c (注册协议)
            │ ip_rcv / arp_rcv  ←── packet_type ───  dev_queue_xmit
 net/ethernet  eth.c (eth_header / eth_type_trans)
 net/core   skbuff.c  dev.c (net_device, netif_rx, backlog, notifier, rtnl)
            │ ndo_start_xmit / ndo_poll
 ═══════════ struct net_device "wlan0" ═════════════════════════════════
 net/wpa    supplicant.c (wpa_supplicant 的活)  eapol.c (四次握手)
            │ cfg80211_scan / cfg80211_connect / control port / add_key
 net/wireless (cfg80211)  core.c  scan.c (BSS 表)  sme.c (连接状态)  util.c
            │ struct cfg80211_ops
 net/mac80211  cfg.c  mlme.c (auth/assoc/deauth)  rx.c  tx.c  key.c  scan.c  main.c
            │ struct ieee80211_ops            ↑ ieee80211_rx_irqsafe / tx_status
 ═══════════ drivers/net/wireless/xradio ═══════════════════════════════
            sta.c (ops)  txrx.c  scan.c  main.c (probe/注册)
            wsm.c (WSM 命令/confirm/indication)
            bh.c  (HIF 队列收发)   fwio.c (固件下载)   hwio.c (HIF 寄存器)
            sdio.c (sdio_driver)
            │ sdio_memcpy_toio / fromio (CMD53)
 drivers/mmc  core.c (SDIO core)  host/sunxi-smhc.c
\end{txtcode}

依赖只向下：驱动不知道 IP，IP 不知道 802.11，supplicant 只通过 cfg80211 接口操作无线设备。每一层的头文件放在 \file{src/include/net/} 下，名字与 Linux 一致（\code{skbuff.h}、\code{netdevice.h}、\code{cfg80211.h}、\code{mac80211.h} $\ldots$），驱动私有头文件留在驱动目录里（\code{xradio.h}、\code{wsm.h}、\code{hwio.h}），和 Linux 的 \file{drivers/net/wireless/xradio/} 一样。

\section{目录与 Linux 对照}

\begin{table}[H]
\centering
\begin{tabular}{p{4.6cm}p{3.4cm}p{5.2cm}}
\toprule
miniOS & Linux 对应 & 作用 \\
\midrule
\code{src/net/core/skbuff.c} & \code{net/core/skbuff.c} & 固定池（32 个）的 \code{sk\_buff}；\code{skb\_push/pull/put} 让每层原地加减头部 \\
\code{src/net/core/dev.c} & \code{net/core/dev.c} & \code{register\_netdev}、notifier、\code{netif\_rx} backlog、\code{dev\_queue\_xmit}、\code{rtnl\_lock}、轮询入口 \code{net\_poll} \\
\code{src/net/ethernet/eth.c} & \code{net/ethernet/eth.c} & 以太网头的构造与解析 \\
\code{src/net/ipv4/*} & \code{net/ipv4/*} & ARP、IPv4 收发、ICMP echo、UDP socket、地址配置、DHCP（\code{ipconfig.c} 对应内核 IP autoconfig） \\
\code{src/net/wireless/*} & \code{net/wireless/}（cfg80211） & wiphy、BSS 表、SME 连接状态、EAPOL control port、802.11$\leftrightarrow$802.3 转换 \\
\code{src/net/mac80211/*} & \code{net/mac80211/} & SoftMAC：组 auth/assoc/deauth 帧、数据帧封装、密钥、扫描、RX tasklet \\
\code{src/net/wpa/*} & 用户态 \code{wpa\_supplicant} & 读 \file{/boot/wifi.txt}、选 AP、RSN IE、PBKDF2、四次握手、装钥、开端口、DHCP \\
\code{.../crypto/wpa\_crypto.c} & \code{lib/crypto/} & PBKDF2/PRF/HMAC-SHA1/AES unwrap \\
\code{src/net/tftp/*} & 无（应用协议） & TFTP 客户端 \\
\code{.../xradio/sdio.c} & \code{xradio/sdio.c} & \code{xradio\_sdio\_driver}：probe、CMD53 自检、提速 \\
\code{.../hwio.c} & \code{xradio/hwio.c} & HIF 寄存器、APB/AHB 间接访问 \\
\code{.../fwio.c} & \code{xradio/fwio.c} & DPLL、bootloader、固件 FIFO 下载、SDD \\
\code{.../bh.c} & \code{xradio/bh.c} & 底半部：HIF 队列读写、序号、固件输入缓冲计数 \\
\code{.../wsm.c} & \code{xradio/wsm.c} & WSM 请求编码、confirm 等待、indication 分派 \\
\code{.../main.c} & \code{xradio/main.c} & 核心 probe：装固件 $\rightarrow$ WSM 启动 $\rightarrow$ CONFIGURATION $\rightarrow$ 注册到 mac80211 \\
\code{.../sta.c} & \code{xradio/sta.c} & \code{ieee80211\_ops}：start、JOIN 与关联参数、set\_key、poll \\
\code{.../scan.c} & \code{xradio/scan.c} & \code{hw\_scan} $\rightarrow$ WSM START\_SCAN；完成通知 \\
\code{.../txrx.c} & \code{xradio/txrx.c} & 802.11 帧 $\leftrightarrow$ WSM TX/RX；剥 IV/MIC；TX confirm \\
\bottomrule
\end{tabular}
\end{table}

\section{接口边界}

\textbf{驱动 $\leftrightarrow$ mac80211}（\file{include/net/mac80211.h}）：驱动实现 \code{ieee80211\_ops}（\code{start}、\code{tx}、\code{hw\_scan}、\code{bss\_info\_changed}、\code{set\_key}、\code{poll}），在 probe 里 \code{ieee80211\_alloc\_hw()} + \code{ieee80211\_register\_hw()}。驱动向上只用四个函数：\code{ieee80211\_rx\_irqsafe()}、\code{ieee80211\_tx\_status()}、\code{ieee80211\_scan\_completed()}、\code{ieee80211\_beacon\_loss()}。驱动只搬帧和下命令，不构造任何 802.11 管理帧。

\textbf{mac80211 $\leftrightarrow$ cfg80211}（\file{include/net/cfg80211.h}）：mac80211 实现 \code{cfg80211\_ops}（\code{scan}、\code{auth}、\code{assoc}、\code{deauth}、\code{add\_key}、\code{tx\_control\_port}）。cfg80211 保存扫描得到的 BSS 表，\code{cfg80211\_connect()} 依次调 auth、assoc。

\textbf{cfg80211 $\leftrightarrow$ supplicant}：supplicant 调 \code{cfg80211\_scan()}、\code{cfg80211\_get\_bss()}、\code{cfg80211\_connect()}、\code{cfg80211\_add\_key()}，EAPOL 帧经 control port 收发（Linux 里是 nl80211），握手完成后调 \code{cfg80211\_port\_authorized()} 打开数据端口、置 carrier。

\textbf{net\_device}：mac80211 注册 \code{wlan0}，\code{ndo\_start\_xmit} 把以太网帧转成 802.11 数据帧，\code{ndo\_poll} 驱动整个无线栈。IP 层只看到一个以太网设备。

\section{一个包怎么走}

发送（\code{ping 192.168.1.1}）：

\begin{txtcode}[basicstyle=\ttfamily\scriptsize]
icmp_ping → ip_queue_xmit (加 IP 头, arp_resolve 下一跳)
  → eth_header (加以太网头) → dev_queue_xmit
  → ieee80211_subif_start_xmit (去掉以太网头, 加 LLC/SNAP、CCMP 预留、802.11 头 ToDS)
  → ieee80211_tx_skb (cookie = WSM packet id) → xradio_tx
  → wsm_tx_frame (24 字节 wsm_tx 头, 队列/速率策略) → xradio_bh_tx → CMD53
\end{txtcode}

接收：

\begin{txtcode}[basicstyle=\ttfamily\scriptsize]
xradio_bh_poll → wsm_handle_rx → xradio_rx_cb
  (WSM RX indication → skb, 剥 CCMP IV/MIC, 填 rx_status)
  → ieee80211_rx_irqsafe  [只入队]
mac80211 RX tasklet (ieee80211_poll 在驱动 poll 之后跑):
  管理帧 → 扫描中: cfg80211_inform_bss_frame；已加入: mlme
  数据帧 → ieee80211_data_to_8023 → EAPOL: control port → supplicant
                                   → 其他: eth_type_trans → netif_rx
net_rx_action: backlog → ip_rcv / arp_rcv → icmp_rcv / udp_rcv → socket 队列
\end{txtcode}

一个 skb 从 UDP payload 到 802.11 帧都是原地 \code{skb\_push}，不复制；\code{NET\_SKB\_PAD} 128 字节的头部余量足够放 UDP+IP+以太网/802.11+CCMP+LLC。

\section{连接流程}

\code{wlan0} 注册时 supplicant 收到 \code{NETDEV\_REGISTER} 通知，自动连接；\code{wifi} 命令走的是同一个 \code{wpa\_supplicant\_connect()}。

\begin{txtcode}[basicstyle=\ttfamily\scriptsize]
wpa: 读 /boot/wifi.txt
cfg80211_scan (信道 1-13 被动扫描) → mac80211 hw_scan → xradio scan.c → WSM START_SCAN
   beacon → BSS 表；SCAN_COMPLETE → ieee80211_scan_completed
wpa: 选 SSID，校验 RSN（CCMP 单播，CCMP/TKIP 组播，PSK），PBKDF2 算 PMK
cfg80211_connect
   → ieee80211_mgd_auth: bss_info_changed(BSSID) → xradio: 预加入 MIB + WSM JOIN + 过滤 MIB
                          发 Open System auth，等 AP 回复
   → ieee80211_mgd_assoc: 发 assoc req（带 RSN IE），等 assoc resp
                          bss_info_changed(ASSOC) → xradio: 关联模式/BSS 参数/PM/slot MIB
EAPOL M1..M4 (control port) → eapol.c
   M4 的 TX status 回来后装钥：cfg80211_add_key → set_key → WSM ADD_KEY
cfg80211_port_authorized → carrier on → "wlan0: link up"
ipconfig_dhcp 或 wifi.txt 里的静态地址
\end{txtcode}

\section{执行上下文（没有中断时的规则）}

miniOS 没有接 SDIO 中断，全部轮询，但仍按 Linux 的上下文划分来写：

\begin{itemize}
  \item \textbf{底半部}：\code{xradio\_bh\_poll()} 读 HIF 队列并分派。它调用的回调（\code{ieee80211\_rx\_irqsafe}、\code{ieee80211\_tx\_status}、\code{ieee80211\_scan\_completed}、\code{ieee80211\_beacon\_loss}）只能入队或记录状态，不能发帧、不能下 WSM 命令。bh 用静态接收缓冲，\code{in\_rx} 守卫禁止重入。
  \item \textbf{tasklet}：\code{ieee80211\_poll()} 在驱动 poll 之后运行 mac80211 RX tasklet，这里可以发帧、下命令（例如装钥匙）。\code{in\_tasklet} 防止嵌套。
  \item \textbf{工作项}：\code{net\_schedule\_work()} 相当于 workqueue，在 \code{\_\_net\_poll()} 末尾执行；supplicant 用它在 M4 发送完成后装钥。
  \item \textbf{rtnl\_lock}：一把不可递归的大锁。对外入口（\code{udp\_*}、\code{icmp\_ping}、\code{wpa\_supplicant\_connect}、\code{net\_poll}）自己加锁；带 \code{\_\_} 前缀的函数要求调用者已持锁。\code{ipconfig\_dhcp} 走加锁的 UDP 接口，所以不能在持锁时调用。
  \item 在 \code{net\_rx\_action} 里需要 ARP 时，\code{arp\_resolve} 只发请求并返回 \code{-EBUSY}，不在接收路径里轮询。
\end{itemize}

\section{Kconfig}

\begin{txtcode}
Networking support            src/net/Kconfig
  NET                         core + ethernet + ipv4 + DHCP
  TFTP_CLIENT
  CFG80211                    src/net/wireless/Kconfig
  MAC80211                    src/net/mac80211/Kconfig
  WPA2_PSK                    src/net/wpa/Kconfig（依赖 CFG80211）
Device Drivers
  Network device support      src/drivers/net/Kconfig
    WLAN
      WLAN_XRADIO             依赖 WLAN && SDIO && MAC80211
        WLAN_XRADIO_XR829
\end{txtcode}

Makefile 按同样的开关挑目标文件；树莓派（kernel8.img）构建排除整个 \file{src/net/} 和 \file{src/lib/crypto/}。

\section{加一块新网卡要做什么}

\begin{enumerate}
  \item 在 \file{src/drivers/net/wireless/<厂商>/} 实现 \code{ieee80211\_ops}；
  \item probe 里 \code{ieee80211\_alloc\_hw()}、\code{ieee80211\_register\_hw(hw, mac, rates, n)}；
  \item 底半部收帧时构造 skb、填 \code{IEEE80211\_SKB\_RXCB}，调 \code{ieee80211\_rx\_irqsafe()}；TX 完成调 \code{ieee80211\_tx\_status()}（带上 cookie 与 status）。
\end{enumerate}

有线网卡则直接 \code{register\_netdev()} 一个 \code{net\_device}，实现 \code{ndo\_start\_xmit}/\code{ndo\_poll}，收包时 \code{eth\_type\_trans()} + \code{netif\_rx()}，IP 层和 tftp 不用改。

\section{从旧结构的对应关系}

\begin{table}[H]
\centering
\begin{tabular}{p{6.2cm}p{8.4cm}}
\toprule
旧文件 & 现在 \\
\midrule
\code{src/net/net.c} & \code{net/core/*}、\code{net/ethernet/eth.c}、\code{net/ipv4/*} \\
\code{src/drivers/net/wireless/xradio/sta.c}（扫描、JOIN、auth/assoc、EAPOL、数据帧转换） & 802.11 状态机 $\rightarrow$ \code{net/mac80211/}；BSS 表与连接 $\rightarrow$ \code{net/wireless/}；wifi.txt/RSN/握手/装钥 $\rightarrow$ \code{net/wpa/supplicant.c}；固件 MIB 设置 $\rightarrow$ \code{xradio/sta.c} \\
\code{.../xradio/xr829\_core.c} & \code{sdio.c}、\code{hwio.c}、\code{fwio.c}、\code{main.c} \\
\code{.../xradio/wsm.c} & 传输 $\rightarrow$ \code{bh.c}；命令编码与分派 $\rightarrow$ \code{wsm.c} \\
\code{src/net/eapol.c} & \code{src/net/wpa/eapol.c} \\
\code{src/net/tftp.c}、\code{tftp\_udp.c} & \code{src/net/tftp/} \\
\code{src/lib/wpa\_crypto.c}、\code{include/net/wpa\_crypto.h} & \code{src/lib/crypto/}、\code{include/crypto/} \\
\code{include/net/net.h}、\code{include/net/xradio/*} & \code{include/net/} 下的 \code{skbuff}、\code{netdevice}、\code{if\_ether}、\code{ip}、\code{arp}、\code{icmp}、\code{udp}、\code{ipconfig}、\code{ieee80211}、\code{cfg80211}、\code{mac80211}、\code{wpa} 等头文件；WSM 头移到驱动目录 \\
\code{xr829\_driver} & \code{xradio\_sdio\_driver} \\
\code{xradio\_sta\_connect()} & \code{wpa\_supplicant\_connect(dev)} \\
\code{net\_dhcp} / \code{net\_ping} / \code{net\_set\_addr} / \code{net\_parse\_ip} & \code{ipconfig\_dhcp} / \code{icmp\_ping} / \code{inet\_set\_addr} / \code{in4\_pton} \\
\bottomrule
\end{tabular}
\end{table}

固件交互的顺序和参数（预加入 MIB、JOIN、过滤 MIB、关联 MIB、钥匙槽位 0/1、速率策略）与旧实现逐项保持一致。

\section{验证}

\begin{itemize}
  \item \code{make t113}（\code{-Wall -Wextra}）零警告，内核链接通过。
  \item 主机端联调：把 \file{src/net}、\file{lib/crypto} 和 xradio 的 \code{wsm/main/sta/scan/txrx} 在 x86 上用 ASan/UBSan 编译，\code{bh.c} 换成假固件，另一端是 Python 写的 WPA2 AP。跑通：扫描 $\rightarrow$ JOIN $\rightarrow$ 认证 $\rightarrow$ 关联 $\rightarrow$ 四次握手（双方 MIC 校验、装进固件的 TK/GTK 与 AP 一致）$\rightarrow$ DHCP $\rightarrow$ ping 网关 $\rightarrow$ TFTP 下载 1000 字节并校验 $\rightarrow$ 组播密钥更新 $\rightarrow$ AP 踢下线（链路断开）$\rightarrow$ 重新连接 $\rightarrow$ ping。
  \item 尚未上板：SDIO、固件真实行为需要在 T113 板子上确认。
\end{itemize}

上板时预期的串口日志（顺序）：

\begin{txtcode}
mmc1:0001:1: firmware running
wsm: firmware startup: ...
mmc1:0001:1: configured MAC 02:0a:9e:22:82:01 ...
cfg80211: phy0 registered, 12 bitrates
net: wlan0 registered, hwaddr 02:0a:9e:22:82:01
wpa: /boot/wifi.txt ssid="TP-Link_B114" psk=<set>
cfg80211: passive scan, 13 channels
cfg80211: AP ssid="..." bssid=... channel=... rcpi=... secured
wpa: selected SSID="TP-Link_B114" ...
wlan0: joining BSS
wlan0: open-system authentication
wlan0: association request
wlan0: associated, AID ...
wpa: WPA2 4-way handshake
wpa: WPA2 keys installed
wlan0: link up
dhcp: bound 192.168.x.x gateway ...
\end{txtcode}
